\chapter{Esercizi}
\section{Esercizi VAN e VA}

\subsection{Richiamo delle formule}
\subsection*{Formule utilizzate}

\noindent\textbf{1. Valore attuale di più flussi di cassa}

\[
VA=\sum_{n=1}^{N}\frac{F_n}{(1+i)^n}
\]

\noindent dove:
\begin{itemize}
    \item $F_n$ = flusso di cassa ricevuto nell'anno $n$
    \item $i$ = tasso di attualizzazione
    \item $n$ = anno in cui si verifica il flusso
\end{itemize}

\noindent\textbf{2. Valore Attuale Netto (VAN)}

\[
VAN=F_0+VA
\]

\noindent oppure, in forma completa:

\[
VAN=F_0+\sum_{n=1}^{N}\frac{F_n}{(1+i)^n}
\]

\noindent dove $F_0$ rappresenta l'investimento iniziale ed è generalmente
negativo perché costituisce un'uscita di denaro.

\noindent\textbf{3. Valore attuale di una rendita annua costante}

\[
VA=C\left[
\frac{1}{r}
-\frac{1}{r(1+r)^t}
\right]
\]

\noindent dove:
\begin{itemize}
    \item $C$ = flusso di cassa costante ricevuto ogni periodo
    \item $r$ = tasso di attualizzazione
    \item $t$ = numero di periodi
\end{itemize}

\noindent Questa formula è una scorciatoia utilizzabile quando si riceve
\textbf{lo stesso flusso di cassa per un numero determinato di periodi}.
Una volta trovato il valore attuale della rendita, il VAN si calcola come:

\[
VAN=-\text{Investimento iniziale}+VA
\]
\newpage
\subsection{Esercizio 1}

La società \textbf{InnovaTech S.r.l.} sta considerando di investire in un progetto di espansione aziendale, caratterizzato dai seguenti flussi di cassa:

\begin{itemize}
    \item \textbf{Investimento iniziale:} 100.000 euro.
\end{itemize}

\noindent I flussi di cassa netti previsti per gli anni successivi sono:

\begin{itemize}
    \item Anno 1: 30.000 euro;
    \item Anno 2: 40.000 euro;
    \item Anno 3: 50.000 euro;
    \item Anno 4: 60.000 euro;
    \item Anno 5: 70.000 euro.
\end{itemize}

\noindent Il tasso di sconto appropriato per questo tipo di investimento è del \textbf{10\% annuo}.

\medskip

\noindent Si richiede di:

\begin{enumerate}
    \item decidere se effettuare l'investimento utilizzando il \textbf{VAN};
    \item determinare come cambierebbe la redditività del progetto se il tasso di sconto aumentasse al \textbf{12\% annuo}.
\end{enumerate}

\subsubsection*{Svolgimento}

\noindent\textbf{1. Calcolo del VAN con $i=10\%$}

\noindent Attualizziamo tutti i flussi di cassa previsti nei cinque anni:

\[
VA=
\frac{30.000}{1,10^1}+
\frac{40.000}{1,10^2}+
\frac{50.000}{1,10^3}+
\frac{60.000}{1,10^4}+
\frac{70.000}{1,10^5}
=182.341,6
\]

\noindent Calcoliamo il VAN considerando l'investimento iniziale di 100.000 euro:

\[
VAN=-100.000+182.341,6=82.341,6
\]

\[
\boxed{VAN\simeq82.342\text{ euro}}
\]

\noindent Poiché il $VAN>0$, il progetto \textbf{crea valore} e l'investimento risulta conveniente.

\medskip
\noindent\textbf{2. Calcolo del VAN con $i=12\%$}

\noindent Aumentiamo il tasso di attualizzazione al 12\% e ricalcoliamo il valore attuale:

\[
VA=
\frac{30.000}{1,12^1}+
\frac{40.000}{1,12^2}+
\frac{50.000}{1,12^3}+
\frac{60.000}{1,12^4}+
\frac{70.000}{1,12^5}
=172.113,4
\]

\noindent Il nuovo VAN è quindi:

\[
VAN=-100.000+172.113,4=72.113,4
\]

\[
\boxed{VAN\simeq72.113\text{ euro}}
\]

\noindent Il VAN rimane positivo, quindi il progetto è ancora \textbf{conveniente}. Tuttavia, aumentando il tasso di attualizzazione dal 10\% al 12\%, il valore attuale dei flussi futuri diminuisce e, 
di conseguenza, diminuisce anche il VAN:

\[
i\uparrow\quad\Rightarrow\quad VA\downarrow\quad\Rightarrow\quad VAN\downarrow
\]
\newpage
\subsection{Esercizio 2}

La società \textbf{Alma Chimica S.p.A.} sta valutando l'acquisto di un nuovo macchinario del costo di \textbf{6,5 milioni di euro}.

\noindent I ricavi previsti ammontano a \textbf{3 milioni di euro all'anno}, mentre i costi di esercizio sono pari a \textbf{2 milioni di euro all'anno}.

\noindent Inoltre:
\begin{itemize}
    \item dopo il quinto anno è necessaria una manutenzione straordinaria del costo di \textbf{1 milione di euro};
    \item dopo il decimo anno è necessaria un'ulteriore manutenzione straordinaria del costo di \textbf{1 milione di euro};
    \item al termine del quindicesimo anno il macchinario può essere venduto per \textbf{800.000 euro}.
\end{itemize}

\noindent Considerando un tasso di interesse del \textbf{7\%}, determinare il \textbf{VAN del macchinario}.

\subsubsection*{Svolgimento}

\noindent Il costo iniziale del macchinario è:

\[
F_0=-6.500.000
\]

\noindent Ogni anno il macchinario genera ricavi per 3 milioni di euro e costi di esercizio per 2 milioni di euro. Il \textbf{flusso di cassa netto annuo} è quindi:

\[
F=3.000.000-2.000.000=1.000.000
\]

\noindent Poiché il flusso di 1 milione di euro si ripete ogni anno per 15 anni, possiamo utilizzare la formula del valore attuale di una \textbf{rendita annua costante}:

\[
VA=C\left[\frac{1}{r}-\frac{1}{r(1+r)^t}\right]
\]

\noindent con:
\begin{itemize}
    \item $C=1.000.000$ $\rightarrow$ flusso di cassa netto annuo;
    \item $r=0,07$ $\rightarrow$ tasso di attualizzazione;
    \item $t=15$ $\rightarrow$ durata dell'investimento.
\end{itemize}

\noindent Sostituendo i valori:

\[
VA=
1.000.000
\left[
\frac{1}{0,07}
-
\frac{1}{0,07(1,07)^{15}}
\right]
=9.107.914
\]

\noindent Dobbiamo poi considerare le due \textbf{manutenzioni straordinarie} da 1 milione di euro, sostenute rispettivamente alla fine del quinto e del decimo anno.

\noindent Essendo delle uscite di denaro, vengono considerate con segno negativo e attualizzate separatamente:

\[
VA_{\text{manutenzioni}}
=
-\frac{1.000.000}{(1,07)^5}
-\frac{1.000.000}{(1,07)^{10}}
\]

\[
VA_{\text{manutenzioni}}
=-1.221.335
\]

\noindent Quindi, fino a questo punto, abbiamo:
\begin{itemize}
    \item $+9.107.914$ euro $\rightarrow$ valore attuale dei flussi di cassa annuali;
    \item $-1.221.335$ euro $\rightarrow$ valore attuale delle manutenzioni straordinarie.
\end{itemize}
\noindent Al termine dei 15 anni il macchinario viene venduto per 800.000 euro.
Poiché questo incasso avverrà al quindicesimo anno, dobbiamo attualizzarlo per
conoscere quanto vale oggi:

\[
VA_{\text{vendita}}
=
\frac{800.000}{(1,07)^{15}}
=
289.956,8
\]

\noindent A questo punto possiamo calcolare il \textbf{VAN dell'intero progetto},
sommando tutti i flussi di cassa attualizzati e considerando il costo iniziale
del macchinario:

\[
VAN
=
-6.500.000
+9.107.914
-1.221.335
+289.956,8
\]

\[
\boxed{VAN=1.676.535,80\text{ euro}}
\]

\noindent Il VAN è \textbf{positivo}, quindi il progetto è conveniente:
considerando anche il valore temporale del denaro, l'investimento genera un
valore aggiuntivo di circa 1,68 milioni di euro.
\newpage
\subsection{Esercizio 3}

\noindent La società Digital Technologies S.r.l. sta valutando due progetti di
investimento alternativi, il \textbf{Progetto X} e il \textbf{Progetto Y},
entrambi con durata di 4 anni e caratterizzati dai seguenti flussi di cassa.

\medskip
\noindent\textbf{Progetto X:}
\begin{itemize}
    \item Investimento iniziale: 600.000 euro
    \item Anno 1: 150.000 euro
    \item Anno 2: 180.000 euro
    \item Anno 3: 210.000 euro
    \item Anno 4: 240.000 euro
\end{itemize}

\noindent\textbf{Progetto Y:}
\begin{itemize}
    \item Investimento iniziale: 750.000 euro
    \item Anno 1: 200.000 euro
    \item Anno 2: 220.000 euro
    \item Anno 3: 240.000 euro
    \item Anno 4: 260.000 euro
\end{itemize}

\noindent Considerando un tasso di sconto del 10\% annuo, si determini quale
progetto di investimento finanziare utilizzando il VAN.
\subsubsection*{Svolgimento}

\noindent\textbf{Progetto X}

\noindent Attualizziamo i quattro flussi di cassa utilizzando il tasso di sconto
del 10\%:

\[
VA_X=
\frac{150.000}{1,10^1}
+\frac{180.000}{1,10^2}
+\frac{210.000}{1,10^3}
+\frac{240.000}{1,10^4}
=606.823,30
\]

\noindent Sottraiamo dal valore attuale dei flussi l'investimento iniziale:

\[
VAN_X=-600.000+606.823,30
=\boxed{6.823,30}
\]

\noindent\textbf{Progetto Y}

\noindent Allo stesso modo attualizziamo i flussi di cassa del secondo progetto:

\[
VA_Y=
\frac{200.000}{1,10^1}
+\frac{220.000}{1,10^2}
+\frac{240.000}{1,10^3}
+\frac{260.000}{1,10^4}
=721.535,41
\]

\noindent Sottraiamo l'investimento iniziale:

\[
VAN_Y=-750.000+721.535,41
=\boxed{-28.464,59}
\]

\noindent Il \textbf{Progetto X è da preferire}, perché presenta un VAN positivo
pari a 6.823,30 euro e quindi genera valore per l'impresa.

\noindent Il Progetto Y presenta invece un VAN negativo: i flussi di cassa
attualizzati non sono sufficienti a recuperare l'investimento iniziale.
\newpage
\subsection{Esercizio 4}

\noindent La società InfoTech S.r.l. sta prendendo in considerazione due progetti
di investimento alternativi, il \textbf{Progetto A} e il \textbf{Progetto B},
entrambi di durata quinquennale e caratterizzati dai seguenti flussi di cassa:

\noindent\textbf{Progetto A:}
\begin{itemize}
    \item Investimento iniziale: 110.000 euro
    \item Flusso di cassa netto annuo: 50.000 euro
\end{itemize}

\noindent\textbf{Progetto B:}
\begin{itemize}
    \item Investimento iniziale: 150.000 euro
    \item Flusso di cassa netto annuo: 60.000 euro
\end{itemize}

\noindent Considerando un tasso di sconto del 10\% annuo, si determini quale
progetto di investimento finanziare utilizzando il VAN.
\subsubsection*{Svolgimento}

\noindent Poiché entrambi i progetti generano un flusso di cassa costante per
5 anni, utilizziamo la formula del valore attuale di una rendita annua costante.

\noindent\textbf{Progetto A}

\[
VA_A=
50.000
\left[
\frac{1}{0,10}
-\frac{1}{0,10(1,10)^5}
\right]
=189.539
\]

\noindent Calcoliamo quindi il VAN sottraendo l'investimento iniziale:

\[
VAN_A=-110.000+189.539
=\boxed{79.539}
\]

\noindent\textbf{Progetto B}

\[
VA_B=
60.000
\left[
\frac{1}{0,10}
-\frac{1}{0,10(1,10)^5}
\right]
=227.447
\]

\noindent Calcoliamo il VAN:

\[
VAN_B=-150.000+227.447
=\boxed{77.447}
\]

\noindent Entrambi i progetti hanno un VAN positivo e sono quindi convenienti.
Tuttavia, essendo progetti alternativi, viene scelto il \textbf{Progetto A},
poiché presenta il VAN maggiore:

\[
\boxed{VAN_A=79.539 > VAN_B=77.447}
\]

\noindent Anche se il Progetto B genera flussi di cassa annuali maggiori,
richiede anche un investimento iniziale più elevato. Il VAN tiene conto di
entrambi gli aspetti e mostra che il Progetto A crea leggermente più valore.